\chapter{Step 3 — Evaluating and Analysing Risks}
\label{chap:Step 3 — Evaluating and Analysing Risks}

%---------------------------- SECTION ----------------------------
\section{Where Thinking Becomes Judgement}
\paragraph{The first two steps of the risk assessment process were primarily about discovery — finding risks and understanding who they affect. This third step is different in character. It is where discovery gives way to analysis, and where analysis gives way to judgement.}
\paragraph{Risk evaluation is the process of taking the risks identified in Step 1, enriched by the exposure analysis of Step 2, and developing a clear, defensible, and actionable picture of their significance. It answers two fundamental questions: how likely is this risk to materialise? And how serious would the consequences be if it did? The combination of those two answers produces something genuinely useful — a prioritised understanding of the risk landscape that tells decision-makers where to focus attention, where to invest in controls, and where, conversely, a lighter touch may be appropriate.}
\paragraph{Done well, risk evaluation transforms a list of identified risks into a genuine management tool. Done poorly — or treated as a mechanical scoring exercise — it produces a false sense of analytical rigour that can be more misleading than helpful.}
\paragraph{This chapter explores how to do it well.}

%---------------------------- SECTION ----------------------------
\section{The Building Blocks — Likelihood and Consequence}
\paragraph{Every approach to risk evaluation, from the simplest qualitative matrix to the most sophisticated quantitative model, is built on the same two fundamental dimensions: likelihood and consequence. Understanding each of these — and the relationship between them — is the foundation of good risk evaluation.}

\subsection{Likelihood}
\paragraph{\textbf{Likelihood} refers to the probability that a risk will materialise — that the hazard will translate into an actual event with real consequences. It is sometimes expressed as a frequency — how often something is expected to occur within a defined time period — and sometimes as a probability — the chance that it will occur at all within that period.}
\paragraph{In practice, estimating likelihood is rarely straightforward. For some risks — those with a long history of occurrence and good data — reasonable quantitative estimates are possible. For others — novel risks, emerging threats, or risks that have never materialised in the organisation's own experience — likelihood can only be estimated qualitatively, drawing on judgement, analogy, and external reference points.}
\paragraph{Several factors influence the likelihood of a risk materialising:}
\begin{itemize}
    \bitem{The effectiveness of existing controls}{a risk subject to strong, well-functioning controls is less likely to materialise than the same risk in an uncontrolled environment.}
    \bitem{The frequency of exposure}{a risk encountered repeatedly is more likely to materialise than one encountered rarely, all else being equal.}
    \bitem{The track record of similar organisations}{industry data, regulatory enforcement records, and published case studies can provide useful reference points for estimating likelihood even where internal data is limited.}
    \bitem{The trajectory of the risk}{is the risk increasing, stable, or decreasing? A risk driven by a rapidly evolving threat — a new form of cyber attack, a changing regulatory environment, a deteriorating macroeconomic context — may warrant a higher likelihood assessment than historical data alone would suggest.}
\end{itemize}

\subsection{Consequence}
\paragraph{\textbf{Consequence} refers to the impact that would result if a risk materialised — the harm, loss, or disruption that would follow. Like likelihood, consequence can be assessed quantitatively — in terms of financial loss, number of customers affected, or regulatory penalty — or qualitatively, using descriptive scales that capture the severity of different types of impact.}
\paragraph{For compliance and legal professionals, consequence assessment needs to be genuinely multidimensional — capturing not just the most immediately obvious impact but the full range of potential consequences across different dimensions:}
\begin{itemize}
    \bitem{Financial consequences}{direct financial losses, regulatory fines and penalties, litigation costs, remediation expenses, and the longer-term financial effects of operational disruption or reputational damage.}
    \bitem{Regulatory and legal consequences}{enforcement action, licence conditions, formal censure, restriction of permissions, personal liability for senior individuals, and the costs and demands of regulatory investigation.}
    \bitem{Reputational consequences}{damage to relationships with customers, regulators, investors, and the broader market; adverse media coverage; loss of stakeholder confidence.}
    \bitem{Operational consequences}{disruption to business activities, loss of key people or systems, supply chain failures, and the costs of recovery and remediation.}
    \bitem{Human consequences}{harm to employees, customers, or third parties; physical, psychological, or financial injury; and in the most serious cases, loss of life.}
    \bitem{Strategic consequences}{longer-term effects on the organisation's ability to pursue its objectives, competitive position, or viability as a going concern.}
\end{itemize}
\paragraph{One of the most important principles in consequence assessment is to consider the \textbf{full chain of consequences} — not just the immediate first-order impact, but the secondary and tertiary effects that flow from it. A data breach causes immediate harm to affected individuals — that is the first-order consequence. It also triggers a regulatory investigation, generates adverse media coverage, and causes customers to withdraw their business — those are second-order consequences. The regulatory investigation results in a significant fine and conditions on the organisation's licence — third-order consequences. Each step in this chain deserves consideration.}

%---------------------------- SECTION ----------------------------
\section{Inherent Risk, Residual Risk, and the Role of Controls}
\paragraph{Before exploring the specific tools and techniques of risk evaluation, there is an important conceptual distinction to address: the difference between \textbf{inherent risk} and \textbf{residual risk}.}
\paragraph{\textbf{Inherent risk} is the level of risk that exists in the absence of any controls — the raw exposure before anything is done to manage it. It reflects the fundamental nature of the hazard and the exposure of those affected by it.}
\paragraph{\textbf{Residual risk} is the level of risk that remains after existing controls have been taken into account — the exposure that the organisation actually faces given the management measures already in place.}
\paragraph{The relationship between these two concepts matters enormously for risk evaluation. An organisation that evaluates only inherent risk may conclude that it faces very high levels of risk across many areas — which may be technically true but provides limited practical guidance, because it ignores the controls that are already in place. An organisation that evaluates only residual risk may underestimate its true exposure if those controls are less effective than assumed, or if they fail.}
\paragraph{Best practice is to assess both — understanding the inherent risk level to appreciate the fundamental nature of the exposure, and the residual risk level to understand what the organisation actually faces in practice. The gap between the two is a measure of the value being delivered by existing controls — and a prompt to consider whether those controls are adequate, well-designed, and actually working as intended.}
\paragraph{This also introduces a third concept: \textbf{target risk} — the level of risk the organisation is aiming to achieve through its control framework. Where residual risk exceeds target risk, further controls or enhancements to existing controls are needed. Where residual risk is significantly below target risk, the organisation may be over-investing in controls relative to the risk — a less common problem, but a real one.}

%---------------------------- SECTION ----------------------------
\section{Qualitative Risk Evaluation}
\paragraph{\textbf{Qualitative risk evaluation} uses descriptive scales — words and categories rather than numbers — to assess likelihood and consequence, and to produce an overall risk rating. It is the most widely used approach in compliance and legal risk assessment contexts, for good reasons: it is accessible, flexible, and does not require quantitative data that may not be available.}

\subsection{Likelihood Scales}
\paragraph{A typical qualitative likelihood scale might use five levels, from rare to almost certain, with descriptive anchors that help assessors apply the scale consistently:}
\begin{table}[H]
  \centering
  \renewcommand{\arraystretch}{1.5}
  \begin{tabularx}{\textwidth}{ | >{\raggedright\arraybackslash}p{0.15\textwidth} 
								| >{\raggedright\arraybackslash}p{0.30\textwidth} 
                                | >{\raggedright\arraybackslash}X | }
    \hline
    \textbf{Level} & \textbf{Label} & \textbf{Descriptive Anchor} \\
    \hline
    1 & \textbf{Rare} & Could happen but only in exceptional circumstances; no known occurrence in the sector in recent years\\
    \hline
    2 & \textbf{Unlikely} & Could happen but would be surprising; has occurred occasionally in the sector\\
    \hline
    3 & \textbf{Possible} & Could happen at some point; occurs from time to time in comparable organisations\\
    \hline
    4 & \textbf{Likely} & Will probably happen at some point; occurs regularly in comparable organisations\\
    \hline
    5 & \textbf{Almost certain} & Expected to happen in most circumstances; has occurred or is occurring\\
    \hline
  \end{tabularx}
  \caption{Likelihood Scales}
\end{table}
\paragraph{The specific labels and anchors matter less than the principle behind them: the scale needs to be defined clearly enough that different people applying it independently would reach broadly consistent conclusions. An undefined scale — where "likely" means one thing to one person and something quite different to another — produces results that cannot be meaningfully aggregated or compared.}

\subsection{Consequence Scales}
\paragraph{A qualitative consequence scale similarly uses defined levels, with descriptive anchors across each of the relevant consequence dimensions:}
\begin{table}[H]
  \centering
  \renewcommand{\arraystretch}{1.5}
  \begin{tabularx}{\textwidth}{ | >{\raggedright\arraybackslash}p{0.10\textwidth} 
								| >{\raggedright\arraybackslash}p{0.15\textwidth} 
								| >{\raggedright\arraybackslash}p{0.25\textwidth} 
								| >{\raggedright\arraybackslash}p{0.25\textwidth} 
                                | >{\raggedright\arraybackslash}X | }
    \hline
    \textbf{Level} & \textbf{Label} & \textbf{Financial} & \textbf{Regulatory} & \textbf{Reputational} \\
    \hline
    1 & \textbf{Negligible} & Immaterial financial impact & Minor technical breach; no regulatory action & No significant external attention\\
    \hline
    2 & \textbf{Minor} & Manageable financial loss & Regulatory inquiry; informal action & Limited adverse coverage; quickly resolved\\
    \hline
    3 & \textbf{Moderate} & Significant but absorbable financial loss & Formal regulatory action; financial penalty & Meaningful adverse coverage; stakeholder concern\\
    \hline
    4 & \textbf{Major} & Serious financial loss; affects performance & Significant regulatory sanction; licence conditions & Serious reputational damage; loss of stakeholder confidence\\
    \hline
    5 & \textbf{Catastrophic} & Existential financial threat & Loss of licence; criminal proceedings & Fundamental reputational harm; potential inability to operate\\
    \hline
  \end{tabularx}
  \caption{Likelihood Scales}
\end{table}
\paragraph{The multidimensional nature of the consequence scale is important. A risk that scores 2 on the financial dimension may score 4 on the regulatory dimension — and it is the highest consequence score across all relevant dimensions that should drive the overall consequence rating, not an average.}

\subsection{The Risk Matrix}
\paragraph{The risk matrix — sometimes called a heat map — is the most widely used tool for combining likelihood and consequence assessments into an overall risk rating. In its most common form, it is a grid with likelihood on one axis and consequence on the other, with the intersection of each combination producing a colour-coded risk rating — typically green for low risk, amber for medium risk, and red for high risk.}
\paragraph{The risk matrix is a genuinely useful tool for visualising the risk landscape and communicating risk priorities to a non-specialist audience. A well-constructed heat map immediately conveys which risks are most significant and where management attention is most needed — in a format that is accessible to board members, senior managers, and regulators alike.}
\paragraph{However, the risk matrix also has well-documented limitations that compliance professionals should understand:}
\paragraph{\textbf{It can create false precision} — the apparently objective numerical output of a risk matrix can obscure the fact that the underlying assessments are qualitative judgements, not measurements. A risk rated 4×4 is not twice as significant as a risk rated 4×2 in any mathematically meaningful sense.}
\paragraph{\textbf{It can compress important distinctions} — in a standard 5×5 matrix, risks at very different points in the likelihood-consequence space may end up in the same colour band. A risk that is almost certain to occur with minor consequences and a risk that is rare but potentially catastrophic may both appear in the amber zone — but they call for very different management responses.}
\paragraph{\textbf{It can encourage rate gaming} — when risk ratings are visible to and scrutinised by senior management, there is a natural tendency for those conducting assessments to adjust their ratings towards the middle of the scale, avoiding the extremes that attract the most attention. This can produce a risk register that looks reassuringly uniform but fails to reflect the genuine range of the organisation's risk exposures.}
\paragraph{None of these limitations are reasons to abandon the risk matrix — it remains a valuable and widely recognised tool. They are reasons to use it thoughtfully, to understand what it shows and what it does not, and to supplement it with narrative analysis that captures the nuances a grid cannot.}

%---------------------------- SECTION ----------------------------
\section{Quantitative Risk Evaluation}
\paragraph{\textbf{Quantitative risk evaluation} uses numerical data, statistical analysis, and probabilistic modelling to produce risk assessments expressed in numerical terms — typically as a probability of occurrence combined with an estimated financial or other quantitative impact.}
\paragraph{Quantitative approaches offer significant advantages where reliable data is available: they produce results that are more precise, more comparable, and more amenable to aggregation and portfolio analysis. They also provide a more rigorous basis for decisions about capital allocation, insurance, and the cost-benefit analysis of control investments.}
\paragraph{In financial services, quantitative risk evaluation is well established — it underpins the capital models, stress tests, and value-at-risk calculations that are central to prudential risk management. In other sectors, quantitative approaches are less universally applied, though they are increasingly used for specific risk types where data quality supports them.}
\paragraph{The main limitation of quantitative approaches — and it is a significant one — is their dependence on data. Where good historical data exists, quantitative models can produce genuinely valuable insights. Where data is limited, absent, or unreliable — as is the case for many of the risks most relevant to compliance professionals, including emerging regulatory risks, reputational risks, and novel operational risks — quantitative modelling can produce spuriously precise results that give false confidence.}
\paragraph{The 2008 financial crisis provided a stark illustration of this risk. Financial institutions that had placed excessive confidence in quantitative risk models — models that were sophisticated in their mechanics but built on historical data that did not capture the possibility of the correlated failures that actually occurred — suffered catastrophic losses. The models were not wrong in their mathematics. They were wrong in their assumptions — and the precision of their outputs obscured rather than revealed that fundamental limitation.}
\paragraph{For compliance professionals, the practical implication is straightforward: use quantitative approaches where data quality genuinely supports them, but do not mistake numerical precision for analytical accuracy. A well-reasoned qualitative assessment grounded in genuine understanding of the risk is often more valuable than a quantitative model built on questionable assumptions.}

%---------------------------- SECTION ----------------------------
\section{Common Analytical Tools and Methodologies}
\paragraph{Beyond the risk matrix, a number of more specialised analytical tools and methodologies are used in risk evaluation — particularly for complex, high-stakes, or technically demanding risk assessments. The following are among the most commonly encountered in compliance and legal contexts.}

\subsection{FMEA — Failure Mode and Effects Analysis}
\paragraph{\textbf{Failure Mode and Effects Analysis} is a structured technique for systematically examining each component of a process or system and identifying the ways in which it could fail — and the consequences of each failure mode. Originally developed for engineering applications, FMEA is now widely used in healthcare, financial services, and other sectors for the analysis of complex processes.}
\paragraph{In an FMEA, each identified failure mode is assessed on three dimensions: the severity of its consequences, the likelihood of its occurrence, and the likelihood of its detection before it causes harm. These three scores are multiplied together to produce a \textbf{Risk Priority Number (RPN)} — a composite score that indicates which failure modes most urgently require attention.}
\paragraph{For compliance professionals, FMEA is particularly useful for the analysis of specific high-risk processes — regulatory reporting processes, customer onboarding procedures, transaction monitoring systems — where a detailed, structured analysis of potential failure modes is warranted.}

\subsection{HAZOP — Hazard and Operability Study}
\paragraph{\textbf{HAZOP} is a structured technique developed in the chemical and process industries for the systematic examination of complex processes and systems. It uses a set of guide words — more, less, none, reverse, other than — applied to each element of a process to identify deviations from intended operation and their potential consequences.}
\paragraph{Whilst HAZOP originated in an engineering context, its underlying logic — using systematic prompts to identify ways in which a process could deviate from its intended operation — has broader applicability. Compliance professionals involved in the risk assessment of complex operational processes, technology implementations, or regulatory change programmes may find elements of the HAZOP approach a useful prompt for structured thinking about process risk.}

\subsection{Bowtie Analysis}
\paragraph{\textbf{Bowtie analysis} is a visual risk assessment technique that maps the pathways between a hazard and its potential consequences, making explicit the causes that could lead to a risk event and the controls that sit between causes and consequences.}
\paragraph{The bowtie diagram takes its name from its shape: on the left side of the diagram are the causes or threat sources that could lead to the central event; on the right side are the consequences that could flow from it. The central event — the point at which the risk materialises — sits at the knot of the bowtie. Controls are represented as barriers on both sides — prevention controls on the left, limiting the likelihood of the event occurring, and mitigation controls on the right, limiting the severity of consequences if it does.}
\paragraph{Bowtie analysis is particularly valuable for risks where understanding the causal pathways and the role of specific controls is important — where the question is not just how likely and how severe? but through what mechanisms could this happen, and at what points can we intervene? For significant regulatory risks, major operational risks, and risks involving complex chains of causation, the bowtie approach can generate genuine analytical insight that a risk matrix alone cannot.}

\subsection{Scenario Analysis and Stress Testing}
\paragraph{\textbf{Scenario analysis} involves constructing specific, plausible scenarios of how a risk could materialise and working through their consequences in detail. Rather than producing a single point estimate of likelihood and consequence, scenario analysis explores a range of possible futures — including low-probability, high-impact scenarios that standard evaluation techniques may underweight.}
\paragraph{\textbf{Stress testing} is a form of scenario analysis that focuses specifically on extreme or adverse scenarios — asking what would happen to the organisation if a particularly severe version of a risk materialised, or if multiple risks materialised simultaneously. In financial services, stress testing is a formal regulatory requirement — firms are required to demonstrate that they have considered the impact of severe but plausible adverse scenarios on their financial position and capital adequacy.}
\paragraph{For compliance professionals more broadly, scenario analysis is a valuable technique for thinking about tail risks — the rare but potentially catastrophic events that standard risk matrices may not capture adequately. It is also a useful tool for board-level risk discussions, where the ability to explore specific scenarios in narrative form can be more engaging and more illuminating than a review of risk ratings.}

%---------------------------- SECTION ----------------------------
\section{Risk Appetite — Knowing Where the Lines Are}
\paragraph{No discussion of risk evaluation would be complete without addressing \textbf{risk appetite} — the organisation's articulated position on the level and types of risk it is willing to accept in pursuit of its objectives. Risk appetite is the standard against which the outputs of risk evaluation are measured — the answer to the question: \textit{is this level of risk acceptable, or does it require action?}}
\paragraph{Risk appetite is one of the most discussed and least well-implemented concepts in risk management. Many organisations have risk appetite statements — often crafted as part of a regulatory compliance exercise — that are expressed in terms too vague to be operationally useful. \textit{"We have a low appetite for regulatory risk"} is a statement that very few organisations would disagree with, and that provides very little practical guidance for decision-making.}
\paragraph{Well-designed risk appetite frameworks have several characteristics:}

\subsection{They Are Specific}
\paragraph{Useful risk appetite statements are specific enough to guide decisions. Rather than expressing a general preference for low risk, they define — for each material risk type — what level of risk is acceptable, what level is tolerable subject to active management, and what level is not acceptable under any circumstances.}

\subsection{They Are Differentiated}
\paragraph{Not all risks are equal, and an organisation's appetite should reflect that. Many organisations are willing to accept significant strategic and commercial risk in pursuit of growth — that is, after all, the nature of business — whilst having very low or zero tolerance for certain regulatory, conduct, or reputational risks. A risk appetite framework that applies the same level of tolerance across all risk types is unlikely to reflect the genuine preferences and obligations of the organisation.}

\subsection{They Establish Clear Thresholds}
\paragraph{Effective risk appetite frameworks establish \textbf{thresholds} — quantitative or qualitative triggers that indicate when a risk has moved from acceptable to requiring action, or from tolerable to unacceptable. These thresholds connect risk appetite to the risk monitoring and escalation processes discussed in Chapter~\ref{chap:Step 5 - Recording, Reviewing, and Monitoring}.}

\subsection{They Are Owned and Used}
\paragraph{Risk appetite is only useful if it is genuinely owned by the board and senior management — not merely approved as a document — and if it is actively used in risk evaluation and decision-making rather than filed away between annual governance reviews. The test of a good risk appetite framework is not how well-written it is but how consistently it influences the decisions the organisation actually makes.}

\subsection{They Acknowledge Zero-Tolerance Areas}
\paragraph{Most organisations have areas where the risk appetite is effectively zero — where certain types of conduct, certain regulatory breaches, or certain categories of harm are simply not acceptable regardless of commercial context. Making these zero-tolerance areas explicit is both good governance and good risk culture — it provides clear guidance to people throughout the organisation about the lines that must not be crossed.}

%---------------------------- SECTION ----------------------------
\section{Building the Risk Register}
\paragraph{The output of the evaluation step is, in most organisations, captured in a \textbf{risk register} — the central documentary record of the organisation's identified and evaluated risks. The risk register is not merely a compliance artefact. When properly maintained, it is a living management tool — a structured, accessible record of the risk landscape that informs decision-making at every level of the organisation.}
\paragraph{At this stage in the process, each entry in the risk register should capture:.}
\begin{itemize}
    \bitem{Risk description}{a clear, concise statement of what the risk is.}
    \bitem{Risk category}{mapped to the organisation's risk taxonomy.}
    \bitem{Risk owner}{the individual with primary accountability for managing the risk.}
    \bitem{Affected parties}{drawn from the Step 2 analysis.}
    \bitem{Inherent risk rating}{likelihood and consequence before controls, with narrative justification.}
    \bitem{Existing controls}{a description of the controls currently in place.}
    \bitem{Control effectiveness}{an assessment of how well existing controls are working.}
    \bitem{Residual risk rating}{likelihood and consequence after controls, with narrative justification.}
    \bitem{Target risk rating}{the level of risk the organisation is aiming to achieve.}
    \bitem{Required actions}{any additional controls or enhancements needed to move from residual to target risk.}
    \bitem{Review date}{when the entry will next be formally reviewed.}
\end{itemize}
\paragraph{The quality of a risk register is not measured by its length. A register with fifty well-described, accurately evaluated, and actively managed risks is considerably more valuable than one with five hundred entries of varying quality that nobody reads or acts upon. The discipline of maintaining a genuinely useful risk register — rather than a comprehensive but unwieldy one — is one of the marks of a mature risk management function.}

%---------------------------- SECTION ----------------------------
\section{Communicating Risk Evaluations — Making the Analysis Useful}
\paragraph{Risk evaluation only delivers value if its outputs are communicated effectively to the people who need to act on them. For compliance professionals, this means being able to present risk analysis in forms that are appropriate for different audiences — from detailed technical assessments for risk specialists, to clear and actionable summaries for operational managers, to high-level board reports that convey the most material risks without overwhelming with detail.}
\paragraph{A few principles are worth keeping in mind:}
\paragraph{\textbf{Lead with significance, not comprehensiveness} — the most important thing a risk report can tell a board or senior management committee is what matters most and why. A report that presents fifty risks at equal length provides much less decision-useful information than one that clearly identifies the five most significant risks and explains them in depth.}
\paragraph{\textbf{Use narrative alongside numbers} — risk ratings and matrices are useful summary tools, but they do not tell the full story. Narrative explanation — describing the nature of the risk, the reasoning behind the evaluation, and the context in which it sits — is essential for genuine understanding and genuine engagement.}
\paragraph{\textbf{Be honest about uncertainty} — risk evaluations are judgements, not facts. Being explicit about the degree of confidence behind an assessment — acknowledging where data is limited, where assumptions are significant, or where the assessment could reasonably be higher or lower — builds credibility and supports better decision-making.}
\paragraph{\textbf{Make the so-what clear} — every risk communication should leave its audience in no doubt about what is expected of them. Does this risk require a decision? Does it require resource allocation? Does it require escalation? The risk evaluation step feeds directly into the control and action decisions of Step 4 — and the communication of its outputs should make that connection explicit.}

%---------------------------- SECTION ----------------------------
\newpage
\section{Chapter Summary}
In this chapter we learned about the following key points:
\begin{table}[H]
  \centering
  \renewcommand{\arraystretch}{1.5}
  \begin{tabularx}{\textwidth}{ | >{\raggedright\arraybackslash}p{0.35\textwidth} 
                                | >{\raggedright\arraybackslash}X | }
    \hline
    \textbf{Concept} & \textbf{Key Point} \\
    \hline
    \textbf{Likelihood} & How probable is the risk materialising — assessed quantitatively or qualitatively\\
    \hline
    \textbf{Consequence} & The multidimensional impact if the risk occurs — financial, regulatory, reputational, operational, human\\
    \hline
    \textbf{Inherent risk} & Risk level before controls — the raw exposure\\
    \hline
    \textbf{Residual risk} & Risk level after controls — the actual exposure\\
    \hline
    \textbf{Target risk} & The risk level the organisation is aiming to achieve\\
    \hline
    \textbf{Risk matrix} & A widely used but imperfect tool for visualising and communicating risk levels\\
    \hline
    \textbf{FMEA} & Structured analysis of failure modes in complex processes\\
    \hline
    \textbf{HAZOP} & Systematic examination of process deviations and their consequences\\
    \hline
    \textbf{Bowtie analysis} & Visual mapping of causes, controls, and consequences around a risk event\\
    \hline
    \textbf{Scenario analysis} & Exploring specific plausible futures — especially useful for tail risks\\
    \hline
    \textbf{Risk appetite} & The organisation's articulated position on acceptable risk levels\\
    \hline
    \textbf{Risk register} & The central documentary record of identified, evaluated, and managed risks\\
    \hline
  \end{tabularx}
  \caption{Risk Concepts}
\end{table}

\paragraph{Key takeaways:}
\begin{itemize}
  \item{Risk evaluation combines likelihood and consequence assessments to produce a prioritised picture of the risk landscape.}
  \item{Both inherent and residual risk should be assessed — understanding the gap between them reveals the value of existing controls.}
  \item{Qualitative approaches are appropriate and widely used — but require clearly defined scales to produce consistent and comparable results}
  \item{Quantitative approaches offer precision where data quality supports them — but numerical precision should not be mistaken for analytical accuracy.}
  \item{Risk appetite is the standard against which evaluation outputs are measured — it needs to be specific, differentiated, and genuinely used to be effective.}
  \item{The risk register is a living management tool, not a compliance document — its quality is measured by its usefulness, not its comprehensiveness.}
  \item{Risk evaluation outputs must be communicated effectively — leading with significance, using narrative alongside numbers, and being honest about uncertainty.}
\end{itemize}

\barquote{In Chapter~\ref{chap:Step 4 — Deciding on Controls and Actions}, we move from analysis to action — exploring the full range of options available for responding to evaluated risks, how to choose between them, and what it means to design and implement controls that are genuinely effective rather than merely documented.}